%%%%%%%%%%%%%%%%%%%%%%%%%%%%%%%%%%%%%%%%%
% MNotation — JLA Tool Report
% Journal of Learning Analytics
% Template: JLA_article.cls (v1.0, 15/11/17)
%
% Compile with: pdflatex → bibtex → pdflatex × 2
%%%%%%%%%%%%%%%%%%%%%%%%%%%%%%%%%%%%%%%%%

\documentclass[fleqn,10pt]{JLA_article}

\usepackage[english]{babel}
\setlength{\headheight}{38.33pt}
\addtolength{\topmargin}{-27.33pt}
\usepackage{hyperref}
\hypersetup{
  hidelinks, colorlinks, breaklinks=true,
  urlcolor=color2, citecolor=color1, linkcolor=color1,
  bookmarksopen=false,
  pdftitle={MNotation: A Learning Analytics Tool for Human-AI Collaborative Qualitative Coding},
  pdfauthor={[Authors]}
}
\usepackage{booktabs}
\usepackage{enumitem}
\usepackage{graphicx}
\usepackage{amsmath}

%----------------------------------------------------------------------------------------
%	ARTICLE INFORMATION
%----------------------------------------------------------------------------------------

\PaperTitle{MNotation: A Learning Analytics Tool for Human--AI Collaborative Qualitative Coding in Educational Research}

\Authors{[Author 1]\textsuperscript{1}*, [Author 2]\textsuperscript{1}, [Author 3]\textsuperscript{2}}

\affiliation{\textsuperscript{1}\textit{[Department, University, City, Country. email: author@university.edu}}
\affiliation{\textsuperscript{2}\textit{[Department, University, City, Country. email: author2@university.edu}}

\Keywords{learning analytics, qualitative coding, human--AI collaboration, active learning, thematic analysis, LLM-assisted annotation, trace data}

\Submitted{DD/MM/YYYY}
\Accepted{DD/MM/YYYY}
\Published{DD/MM/YYYY}

\Volume{XX}
\Number{X}
\Pages{1--20}
\Doi{10.18608/jla.YYYY.XXXX}

%----------------------------------------------------------------------------------------
%	NOTES FOR PRACTICE
%----------------------------------------------------------------------------------------

\Notesname{Notes for Practice}

\note{MNotation can be deployed by any research team conducting qualitative coding at scale.
      It requires no local installation, runs in any web browser, and supports dozens of
      simultaneous participants via QR code or URL — making it accessible for both in-person
      workshops and remote or hybrid settings.}

\note{The three-phase design separates independent human judgement from AI-assisted review,
      enabling researchers to reflect on their own coding decisions as a deliberate epistemic
      activity. Annotation times, accept/override choices, and interaction events are all
      captured automatically, producing a process-level dataset alongside the final coded corpus.}

\note{Instructors can use MNotation in research methods courses to make qualitative coding
      tangible and measurable: students' annotation patterns, decision latencies, and
      AI-override behaviour become pedagogically meaningful data that instructors can
      project and discuss in real time.}

\note{The active learning module identifies the most conceptually ambiguous text segments
      and routes them to human annotators for priority review, substantially reducing the
      amount of material requiring expert attention — a practical efficiency gain for teams
      working with large educational corpora.}

%----------------------------------------------------------------------------------------
%	ABSTRACT
%----------------------------------------------------------------------------------------

\Abstract{%
Qualitative coding is foundational to educational research but difficult to scale: it is
time-intensive, requires expert judgement, and generates little process data about how
interpretive decisions are made. Large language models offer speed, but their agreement
with human researchers in nuanced educational contexts is poorly characterised.
This paper introduces \textbf{MNotation}, an open-source learning analytics platform that
structures human--AI collaboration in qualitative coding through a three-phase workflow:
independent human annotation, LLM-assisted review, and active-learning-guided prioritisation
of the most uncertain segments. MNotation captures fine-grained behavioural trace data at
every phase — annotation latencies, AI accept/override decisions, and interaction events —
making the human--AI meaning-negotiation process itself an object of learning analytics inquiry.
We demonstrate MNotation through two live deployments: a Hong Kong research seminar (March 2026,
$n=69$ active participants) and a Beijing workshop (May 2026, $n=82$). Human--LLM agreement
rates of 31.9\% and 16.7\% respectively, and strikingly divergent AI-label override rates
(27.6\% vs.\ 0.2\%), reveal how context shapes human--AI meaning negotiation in qualitative
coding. The contrasting patterns — Hong Kong participants challenged AI suggestions more
frequently while Beijing participants accepted them at a much higher rate — provide
cross-contextual evidence that MNotation's process data captures interpretive dynamics
unavailable through conventional coding software.
MNotation is freely available at \url{https://github.com/XianghuiMeng-1020/active-labeling}.%
}

%----------------------------------------------------------------------------------------

\begin{document}

\flushbottom
\maketitle
\thispagestyle{fancy}

%----------------------------------------------------------------------------------------
\section{Introduction}
%----------------------------------------------------------------------------------------

Qualitative thematic analysis — the systematic process of reading texts, identifying
patterns, and assigning interpretive codes — remains one of the most widely used
methods in educational research \cite{braun2006using}. It is also one of the most
demanding: analysing hundreds of interview transcripts, student essays, or discussion
forum posts requires expert judgement, sustained engagement with text, and careful
negotiation among coders. As educational datasets grow in scale — driven by the
proliferation of online learning platforms, AI-assisted tutoring systems, and
large-scale studies — the mismatch between the richness of qualitative inquiry
and the constraints of human labour has become a pressing methodological challenge
\cite{papamitsiou2014learning}.

Large language models (LLMs) offer an apparent solution: they can label text at
machine speed. Recent work suggests that LLMs can approximate human inter-rater
reliability for well-defined coding schemes \cite{bing2024llm}, but questions
remain about whether automated coding is methodologically defensible for the
boundary-ambiguous, theory-laden constructs that characterise educational research.
More fundamentally, treating the LLM as an automated coder bypasses the epistemic
value of the coding process itself: the sustained engagement with text that produces
not just labels but interpretive insight.

A more productive framing positions the LLM not as a replacement for the human coder
but as a \emph{collaborator} — one whose proposals the human researcher can accept,
contest, or revise. This framing, grounded in accounts of human--machine meaning
negotiation \cite{liu2026meaning}, reconceptualises coding as a dialogue between human
understanding and machine inference. It preserves the interpretive depth of qualitative
work while using AI to accelerate the process and surface disagreements that warrant
closer attention.

From a learning analytics perspective, this human--AI negotiation is itself a
data-rich process \cite{gasevic2015lets}. Every moment a researcher pauses before
accepting an AI suggestion, or reaches for a different label, reveals something about
the conceptual boundaries of the coding scheme. Capturing and analysing these decision
traces can illuminate which text segments are genuinely contested, where the coding
taxonomy may be under-specified, and how annotators' judgements evolve over the session.
Yet few tools exist that are designed both to facilitate this negotiation and to generate
the trace data needed to study it.

This paper introduces \textbf{MNotation} (Meaning Notation), an open-source learning
analytics platform that addresses this gap. MNotation organises qualitative coding into
three sequential phases — independent human annotation, LLM-assisted review, and
active-learning-guided prioritisation — and logs rich behavioural data at each phase.
The resulting dataset makes the coding process transparent and analytically tractable,
enabling researchers not only to produce a labelled corpus but to study the process by
which that corpus came to be.

We report a case study deployment at a research seminar where 69 participants annotated
sentences from student essays about AI literacy. The study documents MNotation's
architecture, its novel active learning algorithm (ED-AL v1), and the patterns of
human--AI agreement and disagreement it revealed. Our goal is to demonstrate MNotation's
value for the learning analytics community as both a practical research instrument and a
new lens on human--AI collaborative sense-making in educational contexts.

%----------------------------------------------------------------------------------------
\section{Theoretical Grounding}
%----------------------------------------------------------------------------------------

\subsection{Qualitative Coding as a Learning Activity}

The act of coding qualitative data is not merely a mechanical classification task.
It requires the coder to hold a conceptual framework in mind, interpret each text
segment in relation to that framework, and make explicit decisions about category
boundaries. Repeated engagement with an annotation scheme constitutes a form of
\emph{schema development} \cite{chi2021translating}: coders progressively refine
their understanding of each category through exposure to edge cases and through
negotiation with co-coders.

This perspective motivates the collection of process-level trace data during coding.
Just as learning analytics researchers instrument digital environments to capture
reading behaviours, help-seeking, and self-regulation \cite{flora2025}, MNotation
instruments the coding interface to capture annotation decisions, response latencies,
and revisions. The resulting data support analyses of how individual and group
understanding of a coding scheme develops over time — extending learning analytics
methods from \emph{student} learning to \emph{researcher} learning.

\subsection{Human--AI Meaning Negotiation}

The concept of meaning saturation \cite{liu2026meaning} proposes that meaning in
qualitative analysis is not fixed in a text but emerges from the encounter between
text, reader, and interpretive community. AI-generated labels are not right or wrong
in any absolute sense; they represent one possible reading that the human coder can
engage with critically. The process of accepting, modifying, or rejecting an AI label
is an act of meaning negotiation that surfaces the coder's own interpretive commitments.

This framing has direct implications for tool design. Rather than presenting AI
predictions as authoritative outputs, MNotation frames them as \emph{proposals} that
invite human evaluation. The interface deliberately separates the human coding phase
from the AI-assisted review, ensuring that participants form their own judgements before
encountering the model's suggestions — mirroring established practices for avoiding
anchoring bias in expert judgement \cite{nickerson1998confirmation}.

\subsection{Active Learning in Qualitative Research}

Active learning — in the machine learning sense — refers to algorithms that query human
annotators strategically, directing attention to the examples most likely to improve
model performance \cite{settles2012active}. In the context of qualitative educational
research, active learning serves an additional function: it identifies the text segments
where human--AI disagreement is highest, which are often the segments most conceptually
interesting to the researcher.

MNotation incorporates an active learning module (ED-AL v1, described in Section 3.3)
that combines uncertainty sampling and diversity-based selection to identify a small
subset of segments for priority annotation. This reflects the growing recognition in
learning analytics that directing human attention to the most consequential examples is
itself a form of pedagogical scaffolding \cite{munshi2023analysing} — applicable here to
the scaffolding of the researcher's own analytical process.

%----------------------------------------------------------------------------------------
\section{Tool Description}
%----------------------------------------------------------------------------------------

\subsection{System Overview}

MNotation is a browser-based platform designed for groups of researchers or students to
annotate text corpora collaboratively, in real time, with AI assistance. It is hosted on
Cloudflare's global edge network, which means it can support dozens of simultaneous users
without local server setup and is accessible from any device with a web browser, including
smartphones. This design reflects a commitment to accessibility: annotation workshops can
be run in seminar rooms, online, or in hybrid settings without requiring participants to
install software or create accounts. Participants join via a shared URL or QR code.

The platform has two interfaces: a \emph{participant interface} guiding users through the
three annotation phases, and an \emph{administrator dashboard} providing real-time
monitoring, data management, and active learning control.

\subsection{Three-Phase Annotation Workflow}

\subsubsection{Phase 1: Independent Human Annotation}

Each participant reads text segments presented one at a time and assigns a label from
the predefined taxonomy by tapping or clicking. No AI information is shown during this
phase, ensuring that initial judgements are uncontaminated by machine suggestions.
The interface records three engagement metrics for each annotation decision:

\begin{itemize}[noitemsep]
  \item \texttt{active\_ms} — time the interface was in active browser focus
  \item \texttt{idle\_ms} — time the interface was visible but the tab was unfocused
  \item \texttt{blur\_count} — number of tab-switch events during the annotation
\end{itemize}

A configurable minimum engagement threshold (default: 800 ms) filters accidental or
inattentive clicks before a label is accepted as valid.

\subsubsection{Phase 2: LLM-Assisted Review}

After completing Phase 1, participants enter a review mode in which each segment is
displayed alongside the LLM's predicted label. They can accept the prediction (one tap),
override it using a slide-up selector, or switch between two pre-configured prompt
conditions — a zero-shot prompt (Prompt 1) and a few-shot exemplar prompt (Prompt 2) —
or enter a custom prompt (rate-limited to five queries per session). The tool logs the
final submitted label, whether it differs from the AI's prediction, and the decision
latency. This phase enables participants to critically evaluate a second perspective on
each text — a process analogous to structured peer review in collaborative learning.

\subsubsection{Phase 3: Active Learning Prioritisation}

A subset of text segments — selected by the ED-AL v1 algorithm (Section 3.3) — is
presented for re-annotation. These are the segments where collective human uncertainty is
highest and content diversity is greatest. This phase directs collective attention to the
hard cases: those where the coding taxonomy's boundaries are most contested and where
additional expert deliberation is most valuable.

\subsection{The ED-AL v1 Active Learning Algorithm}

The active learning module uses two sequential steps to identify which segments to
surface in Phase 3.

\textbf{Step 1 — Uncertainty scoring.} For each candidate segment, the LLM is queried
$n$ times with temperature $\tau > 0$ to sample a label distribution. Shannon entropy
is computed over the resulting label frequencies and normalised to $[0, 1]$:

\begin{equation}
  H_{\text{norm}}(x) = \frac{-\sum_{i} p(l_i \mid x)\log_2 p(l_i \mid x)}{\log_2 |L|}
\end{equation}

\noindent where $|L|$ is the number of labels in the taxonomy. High entropy indicates
genuine ambiguity relative to the taxonomy.

\textbf{Step 2 — Diversity selection.} From the top-$H$ highest-entropy candidates,
TF-IDF vectors are computed and a greedy $k$-centre algorithm selects $M$ segments
maximally spread in the feature space. This ensures that the segments sent to human
annotators cover different conceptual regions rather than clustering around a single
ambiguous sentence type.

All algorithm parameters — pool size, entropy threshold, number of LLM samples,
temperature, and final selection count — are adjustable through the administrator
dashboard without code changes.

\subsection{Administrator Dashboard and Data Export}

The administrator interface provides live visualisations of annotation progress: bar charts
showing label distribution per phase update in real time as participants submit decisions.
A participant progress table displays each user's phase completion status. A \emph{freeze}
function pauses the live display, enabling facilitators to discuss emerging patterns with
participants — turning the annotation activity into a collective learning moment in its
own right.

Data export is available in CSV and JSON formats and includes the fields shown in Table~\ref{tab:schema}.

\begin{table}[hbt]
\caption{Key fields in MNotation's per-annotation export.}
\label{tab:schema}
\centering
\small
\begin{tabular}{lp{5.5cm}}
\toprule
\textbf{Field} & \textbf{Description} \\
\midrule
\texttt{session\_id}           & Anonymised participant session identifier \\
\texttt{unit\_id}              & Text segment identifier \\
\texttt{manual\_label}         & Phase~1 human-assigned label \\
\texttt{llm\_predicted}        & LLM-generated label \\
\texttt{user\_accepted\_label} & Final label after Phase~2 review \\
\texttt{active\_ms}            & Active engagement time (ms) \\
\texttt{idle\_ms}              & Idle/background time (ms) \\
\texttt{blur\_count}           & Tab-switch events during annotation \\
\texttt{is\_valid}             & Whether \texttt{active\_ms} $\ge$ threshold \\
\texttt{llm\_mode}             & Prompt condition (zero-shot / few-shot / custom) \\
\bottomrule
\end{tabular}
\end{table}

\subsection{Technical Architecture}

MNotation's backend is implemented in TypeScript on Cloudflare Workers with Hono 4.x as
the routing framework. Data is stored in Cloudflare D1 (serverless SQLite; 12 tables).
Real-time push notifications are delivered via Cloudflare Durable Objects with
server-sent events (SSE). The frontend is built with React 19 and TypeScript, compiled
with Vite 7. The LLM integration defaults to Qwen-Plus via the DashScope API (OpenAI-compatible),
with configurable fallback to any other OpenAI-compatible provider. The complete codebase,
database schema, and deployment documentation are available at
\url{https://github.com/XianghuiMeng-1020/active-labeling}.

%----------------------------------------------------------------------------------------
\section{Case Studies: Two AI Literacy Annotation Deployments}
%----------------------------------------------------------------------------------------

MNotation was deployed in two live research-seminar settings: a university in Hong Kong
in March 2026, and a university in Beijing in May 2026. Both sessions used an identical
coding task and taxonomy, enabling systematic cross-contextual comparison of human--AI
collaboration dynamics.

\subsection{Shared Coding Task and Taxonomy}

The annotation corpus comprised 15 sentences drawn from student-written essays on the topic
of AI literacy, selected to span a range of thematic categories and to include cases with
genuine category ambiguity. The coding taxonomy comprised six codes:
\textbf{CODE} (definitional content), \textbf{EXPLANATION} (explanations of AI concepts),
\textbf{EVALUATION} (assessments of AI reliability or risk), \textbf{RESPONSIBILITY}
(ethics and accountability), \textbf{APPLICATION} (practical uses of AI literacy), and
\textbf{IMPLICATION} (broader consequences or future significance).
A 10-minute briefing with two worked examples preceded annotation in both deployments.

\subsection{Deployment 1 — Hong Kong (March 24, 2026)}

The Hong Kong session was held at a research seminar attended by graduate students and
faculty from education, linguistics, and computer science. The activity was framed as a
participatory demonstration: participants were positioned as co-investigators rather than
assessed subjects. Participants accessed MNotation by scanning a QR code on the seminar
screen. A total of 103 unique users provided consent; \textbf{69 (67\%)} submitted at
least one label (active participants). Among active participants, 53 (77\%) completed the
15-item manual annotation phase and 27 (39\%) completed the LLM-assisted review.
Qwen-Plus was used as the LLM backend. A post-session survey was administered, yielding
18 complete responses.

\subsection{Deployment 2 — Beijing (May 26, 2026)}

The Beijing session was conducted in the evening (19:04--20:08 Beijing Time) with students
from an education and AI literacy programme. Eighty-six unique users provided consent;
\textbf{82 (95\%)} submitted at least one label. Among active participants, 74 (90\%)
completed Phase~1 and 60 (73\%) completed Phase~2 — substantially higher completion rates
than in Hong Kong. No post-session survey was administered for this cohort.

\subsection{Participation and Annotation Patterns}

Table~\ref{tab:cohorts} summarises participation and annotation metrics across both
deployments. The higher completion rates in Beijing may reflect a more structured
workshop environment in which participants remained engaged for the full session duration.
The dramatic difference in annotation speed — mean active time of 8.3 s per sentence in
Hong Kong versus 4.5 s in Beijing — suggests different deliberation strategies: Hong Kong
participants spent more time per item before committing to a label, while Beijing
participants made decisions more quickly. This speed gap is present across all phases
and persists when controlling for session length.

Both cohorts chose \textbf{EXPLANATION} as the most common label, but the Hong Kong
distribution was markedly skewed toward EXPLANATION (28.1\%) while Beijing showed a
substantially more even spread across all five active categories (18--23\%).
The \textbf{CODE} label was unused by participants in both cohorts despite its inclusion
in the taxonomy, suggesting a shared tendency to interpret definitional sentences as
\emph{explanatory} — a cross-contextual signal of boundary ambiguity that warrants
coding scheme revision.

\begin{table}[hbt]
\caption{Participation and annotation metrics for both deployments.}
\label{tab:cohorts}
\centering
\small
\begin{tabular}{lrr}
\toprule
\textbf{Metric} & \textbf{Hong Kong (Mar 24)} & \textbf{Beijing (May 26)} \\
\midrule
Unique users (consent, window)   & 103    & 86    \\
Active participants ($\ge$1 label) & 69   & 82    \\
Completed Phase~1 (manual)       & 53 (77\%)  & 74 (90\%) \\
Completed Phase~2 (LLM review)   & 27 (39\%)  & 60 (73\%) \\
Valid annotation records         & 702   & 2,011 \\
Mean active annotation time      & 8.3 s & 4.5 s \\
Human--LLM agreement             & 31.9\% & 16.7\% \\
AI label override rate           & 27.6\% & 0.2\% \\
Post-session survey responses    & 18    & --    \\
\bottomrule
\end{tabular}
\end{table}

\begin{table}[hbt]
\caption{Phase~1 human label distribution (valid annotations only) for both deployments.}
\label{tab:labels}
\centering
\begin{tabular}{lrrrr}
\toprule
& \multicolumn{2}{c}{\textbf{Hong Kong}} & \multicolumn{2}{c}{\textbf{Beijing}} \\
\cmidrule(lr){2-3}\cmidrule(lr){4-5}
\textbf{Label} & \textbf{Count} & \textbf{\%} & \textbf{Count} & \textbf{\%} \\
\midrule
EXPLANATION    & 243 & 28.1\% & 259 & 22.6\% \\
EVALUATION     & 193 & 22.3\% & 238 & 20.7\% \\
APPLICATION    & 169 & 19.5\% & 209 & 18.2\% \\
RESPONSIBILITY & 137 & 15.8\% & 226 & 19.7\% \\
IMPLICATION    & 123 & 14.2\% & 215 & 18.7\% \\
CODE           &   0 &  0.0\% &   0 &  0.0\% \\
\midrule
\textbf{Total} & \textbf{865} & & \textbf{1,147} & \\
\bottomrule
\end{tabular}
\end{table}

\subsection{Human--LLM Agreement and Override Behaviour}

Among sentence-level comparisons for which both a Phase~1 human label and an LLM
prediction were available, agreement was \textbf{31.9\%} in Hong Kong (177 of 555
comparisons) and \textbf{16.7\%} in Beijing (153 of 915 comparisons). Although both
rates are notably lower than those reported for simpler binary classification tasks,
the difference between cohorts is itself a substantive finding: Hong Kong participants'
Phase~1 labels aligned more closely with LLM predictions, while Beijing participants
coded more independently of the LLM baseline.

The override pattern reveals a striking reversal. In Hong Kong, \textbf{153 of the 555
LLM labels reviewed in Phase~2 were modified (27.6\% override rate)}, indicating that
participants actively deliberated over AI proposals. In Beijing, only \textbf{2 of 905
accepted labels were overridden (0.2\%)}. Together, these patterns suggest two distinct
human--AI collaboration modes: one characterised by critical engagement (Hong Kong) and
one by rapid acceptance (Beijing). Whether this reflects differences in AI literacy,
workshop framing, or time pressure cannot be determined from the current data alone,
but the contrast is precisely the type of cross-contextual signal that motivates
continued deployment and study.

Agreement varied substantially across the 15 sentences in the Hong Kong cohort
(Table~\ref{tab:unit_agree}). The sentence ``\emph{Machine learning algorithms work by
identifying patterns in data}'' achieved 76\% human--LLM agreement, while
``\emph{Companies must ensure transparency in how they use AI systems}'' achieved only 21\%.
The most common disagreement patterns were:
\textsc{Explanation}$\to$\textsc{Evaluation} (41 cases),
\textsc{Application}$\to$\textsc{Responsibility} (36 cases), and
\textsc{Evaluation}$\to$\textsc{Application} (33 cases) — all clustering around
conceptually adjacent category pairs where the \emph{rhetorical function} of a sentence
is ambiguous.

\begin{table}[hbt]
\caption{Human--LLM agreement rates (Hong Kong cohort) for the five most and five
  least agreed-upon sentences.}
\label{tab:unit_agree}
\centering
\small
\begin{tabular}{p{5.5cm}rll}
\toprule
\textbf{Sentence (truncated)} & \textbf{Agr.} & \textbf{Human} & \textbf{LLM} \\
\midrule
Machine learning algorithms\ldots & 76\% & EXPLANATION & EXPLANATION \\
Citizens can use fact-checking\ldots & 64\% & APPLICATION & APPLICATION \\
Without proper governance\ldots & 58\% & IMPLICATION & IMPLICATION \\
As AI continues to advance\ldots & 50\% & IMPLICATION & IMPLICATION \\
AI literacy refers to\ldots & 32\% & EXPLANATION & EXPLANATION \\
\midrule
Bias in AI systems\ldots & 24\% & RESPONSIBILITY & APPLICATION \\
Companies must ensure transparency\ldots & 21\% & RESPONSIBILITY & APPLICATION \\
Developers have an ethical obligation\ldots & 12\% & RESPONSIBILITY & EVALUATION \\
Healthcare professionals are using\ldots & 12\% & APPLICATION & EXPLANATION \\
Research shows that people with higher\ldots &  9\% & EVALUATION & EXPLANATION \\
\bottomrule
\end{tabular}
\end{table}

\subsection{Participant Perceptions (Hong Kong Cohort)}

Eighteen participants completed the post-session survey. Participants reported strong
understanding of the session's conceptual framing: \emph{``I understood that the
AI-supported system was designed to support, not replace, human interpretation''} received
the highest Likert rating ($M = 4.67/5$). The item \emph{``The AI-supported system helped me
generate initial ideas for coding''} received $M = 4.00$. However, \emph{``I was sometimes
unsure why certain excerpts were selected by the system''} also received notable agreement
($M = 3.50$), revealing a \emph{transparency gap} in the active learning selection rationale.

Open-ended responses further characterised the collaborative dynamic. Participants described
the AI as useful for transforming a generative decision into a critical one: ``give a code so
that humans can do judgement [tasks] rather than selection tasks.'' One participant
characterised the experience as having ``a peer to support me in evaluating and rethinking
my initial coding.'' Constructive feedback called for greater transparency about the
system's selection logic and additional task time.

%----------------------------------------------------------------------------------------
\section{Discussion}
%----------------------------------------------------------------------------------------

\subsection{What MNotation Makes Visible}

The two deployments illustrate MNotation's central contribution: it makes the human--AI
coding process \emph{measurable} at a granularity unavailable through conventional
qualitative software. By capturing annotation decisions in sequence — first human, then
AI-assisted, then active-learning-prioritised — the tool produces a longitudinal trace of
how each participant's interpretation intersects with, diverges from, and ultimately
converges with or departs from LLM proposals. This trace is the unit of analysis for a
new genre of learning analytics research: studies of human--AI deliberation in educational
inquiry.

Critically, the comparison across two independent deployments elevates the tool's
evidential value beyond what any single session can provide. The divergence in human--LLM
agreement (31.9\% vs.\ 16.7\%) and override rate (27.6\% vs.\ 0.2\%) between Hong Kong
and Beijing participants, annotating the \emph{same} text with the \emph{same} taxonomy
and \emph{same} LLM backend, demonstrates that these metrics are sensitive to contextual
factors such as workshop structure, participant familiarity with AI systems, and the
temporal dynamics of the session. These rates are not measures of LLM quality; they are
measures of the interpretive space between human researchers and language models in a
given context. That space — and how it shifts — is where the interesting research
questions live.

\subsection{The Active Learning Transparency Gap}

A key finding from the Hong Kong deployment is participants' moderate uncertainty about why
the active learning algorithm selected certain text segments ($M = 3.50$ on the relevant
survey item). This transparency gap is well-documented in human--AI interaction research
\cite{amershi2019software}: users engage more productively with algorithmic recommendations
when they can see the rationale. In MNotation's current implementation, uncertainty scores
and diversity ranks are computed but not surfaced to participants.

The override pattern from Beijing raises a related concern from a different direction.
A 0.2\% override rate — essentially universal acceptance of LLM suggestions in Phase~2 —
could reflect genuine agreement with the AI, rapid task completion under time pressure, or
insufficient critical engagement with the AI's proposals. These possibilities are empirically
distinguishable using the annotation timing data MNotation captures: if override decisions
co-occur with longer active annotation times, they reflect deliberation; if they are absent
even when annotation latency is high, they may indicate a different dynamic. Future studies
should use MNotation's process data to disentangle these patterns.

Future versions should expose uncertainty scores in the Phase~3 interface — showing that a
sentence was flagged because annotators assigned five different labels across responses —
transforming Phase~3 into a collective learning opportunity where participants can reflect
on aggregate disagreement as evidence about the shared intelligibility of the coding scheme.

\subsection{Contributions to Learning Analytics}

MNotation contributes to the learning analytics community in three complementary ways.

\textbf{First, it enables the study of human--AI calibration as a learning process.}
When a researcher sees an AI label that differs from their own, they face a decision
that is simultaneously methodological (which label is correct?) and epistemic (how
confident am I in my reading?). The override rate, annotation latency on revised
decisions, and patterns of multi-phase label change are all indicators of this
calibration process. Analysing these indicators across researchers, tasks, and coding
schemes can produce evidence-based recommendations for AI-assisted qualitative research
practice \cite{jarvela2023human}.

\textbf{Second, it produces datasets that are both analytic outputs and learning
analytics artefacts.} The annotated corpus that emerges from a MNotation session is
useful for downstream analysis; but the session trace data — timing, agreement patterns,
active learning priorities — constitute a learning analytics dataset in their own right,
capturing how a community of researchers collectively negotiated meaning around a shared
corpus. This dual nature positions MNotation at the intersection of tool and instrument
in the learning analytics ecosystem \cite{gasevic2023empowering}.

\textbf{Third, it enables principled scaling of qualitative analysis.} The case study
data illustrate this directly. Among the 15 sentences, human--LLM agreement ranged from
9\% to 76\%. For high-agreement sentences, automated coding may be methodologically
defensible if the researcher documents and reports the agreement evidence. For low-agreement
sentences, human deliberation is clearly warranted. MNotation makes this triage automatic
and evidence-based, enabling research teams to concentrate expert effort where it has the
greatest methodological impact.

%----------------------------------------------------------------------------------------
\section{Limitations}
%----------------------------------------------------------------------------------------

Several limitations should be considered when interpreting findings from both deployments.

\textbf{Sample characteristics.} Participants in both cohorts were graduate students and
faculty at research-intensive universities; their familiarity with qualitative methods and
AI tools likely exceeds that of the broader educational research population. The two
deployments represent different geographic and institutional contexts (Hong Kong and Beijing),
which partially addresses generalisability concerns, but neither cohort was randomly sampled.

\textbf{Single-session design.} Each cohort's data reflect a single session (25--60 minutes),
limiting conclusions about longitudinal learning effects or inter-rater reliability development
over repeated cycles. Future deployments should incorporate multiple sessions with the same
participants to examine how human--AI calibration evolves.

\textbf{Cross-cohort confounds.} The Hong Kong and Beijing sessions differed in time of
day, workshop duration, and degree of facilitator scaffolding, in addition to participant
background. The observed differences in override rate and annotation speed should therefore
be interpreted as hypothesis-generating rather than conclusive evidence of geographic or
cultural differences in human--AI collaboration.

\textbf{LLM specificity.} AI suggestions were generated by Qwen-Plus. Agreement rates and
override patterns may differ with other models (e.g., GPT-4o, Claude, or open-source
alternatives). Researchers deploying MNotation should document and report the model used,
as this is a methodological decision analogous to the choice of inter-rater reliability metric.

\textbf{Survey response rate.} The 18 survey responses represent 26\% of active participants,
limiting the generalisability of perception data. Future work should investigate ways to
increase survey completion within the tool's workflow.

\textbf{Participant verification.} Because MNotation is web-based and accessed via a shared
QR code, some registered sessions reflected researcher testing or participants who rejoined
the session. The analyses reported here use each participant's earliest session and exclude
sessions created before or after the main workshop window. Researchers deploying MNotation
in controlled studies should implement session filtering appropriate to their context.

%----------------------------------------------------------------------------------------
\section{Conclusion and Future Work}
%----------------------------------------------------------------------------------------

MNotation addresses a methodological gap in learning analytics and educational research:
the absence of open, purpose-built platforms for studying how human researchers negotiate
meaning with AI systems during qualitative coding. By separating independent human annotation
from AI-assisted review and active-learning-guided prioritisation, and by logging behavioural
trace data at each phase, the tool creates datasets that are simultaneously useful for research
purposes and informative about the annotation process itself.

The case study reported here demonstrates that the human--AI coding process produces rich,
structured data: the 31.9\% human--LLM agreement rate, the 27.6\% AI override rate,
per-sentence agreement patterns, and annotation timing data each offer a distinct window on
how researchers interpret educational text alongside AI proposals. Together, these data
constitute a new category of learning analytics evidence — not about students learning
subject matter, but about researchers learning to interpret data, with and alongside AI.

Planned developments include: (1) surfacing uncertainty and diversity scores to participants
in Phase~3 to address the transparency gap identified in the case study; (2) support for
open coding (participant-defined categories) in addition to closed taxonomies; (3)
longitudinal session tracking to support inter-session reliability analysis; and (4) a
multi-coder comparison dashboard visualising human--human and human--LLM disagreements
side by side. The tool's modular serverless architecture supports all these extensions
without changes to its core infrastructure.

%----------------------------------------------------------------------------------------
\phantomsection
\section*{Acknowledgments}
\addcontentsline{toc}{section}{Acknowledgments}
%----------------------------------------------------------------------------------------

[Acknowledgments to be added.]

%----------------------------------------------------------------------------------------
\section*{Declaration of Conflicting Interest}
\addcontentsline{toc}{section}{Declaration of Conflicting Interest}
%----------------------------------------------------------------------------------------

The author(s) declared no potential conflicts of interest with respect to the research,
authorship, and/or publication of this article.

%----------------------------------------------------------------------------------------
\section*{Funding}
\addcontentsline{toc}{section}{Funding}
%----------------------------------------------------------------------------------------

[Funding information to be added.]

%----------------------------------------------------------------------------------------
\section*{Tool and Data Availability}
\addcontentsline{toc}{section}{Tool and Data Availability}
%----------------------------------------------------------------------------------------

\textbf{Tool:} \url{https://github.com/XianghuiMeng-1020/active-labeling} (MIT License).
Full deployment documentation, database schema, and API reference are included in the
repository.

\textbf{Datasets:} Two anonymised workshop datasets will be deposited on OSF upon paper
acceptance: (1) the Hong Kong cohort (March 24, 2026: 69 active participants; 702 valid
Phase~1 annotations; 555 human--LLM comparison pairs; 18 survey responses) and (2) the
Beijing cohort (May 26, 2026: 82 active participants; 2,011 valid annotations; 915
human--LLM comparison pairs). Both datasets include session-level participation records,
sentence-level annotation decisions, and fine-grained timing and interaction logs.
Access procedures are documented in the repository README.

%----------------------------------------------------------------------------------------
\phantomsection
\bibliography{mnotation_jla}
%----------------------------------------------------------------------------------------

\end{document}
